\subsection{Genetic Algorithm for Map Generation}

The map generation module uses a genetic algorithm to automatically generate the game map. The aim of this part is not just to place walls randomly, but to create a $30 \times 30$ map that is connected, playable, and suitable for the chase mechanism in this game. In this project, the map structure directly affects the game balance. If the map is too open, monsters may chase the human too easily. If there are too many walls, the movement of both the human and monsters may become restricted, and they may even be separated into unreachable areas. For this reason, the map generation process needs to balance randomness and playability.

The map is represented as a $30 \times 30$ two-dimensional grid. Each cell is stored as a binary value. The value \texttt{0} represents walkable ground, while the value \texttt{1} represents a wall or obstacle. The human and monsters can only move on walkable cells. The map also uses a wrap-around boundary mechanism, which means that when a character moves out of one side of the map, it will appear from the opposite side. Because of this design, modulus operation is used when calculating neighbouring cells, and Breadth-First Search (BFS) can be applied on this wrap-around map structure for connectivity and path distance checking.

In the genetic algorithm, each map is regarded as a chromosome, while a group of maps forms a population. The initial population is generated randomly, where each cell is assigned as either walkable ground or obstacle according to a fixed wall probability. In this implementation, the target wall ratio is set to around $0.28$, which means that about $28\%$ of the map is expected to be walls. This value is used to avoid maps that are either too empty or too crowded.

After the initial maps are generated, each candidate map is evaluated by a fitness function. The fitness function mainly considers four factors: reachability, wall ratio, junction complexity, and path distance. The final fitness score can be expressed as:

\begin{equation}
\text{Fitness} =
S_{\text{reachability}}
+ S_{\text{wall ratio}}
+ S_{\text{junction}}
+ S_{\text{path distance}}
- P_{\text{dead-end}}
\end{equation}

Reachability is used to check whether most walkable cells are connected. Wall ratio controls the proportion of walls and compares it with the target value of $0.28$. Junction complexity measures the number of branches and intersections, which can provide more possible route choices during the chase. Path distance is used to evaluate whether the average BFS path length in the map is suitable for the survival chase setting. The dead-end penalty is added to reduce the number of cells with very few exits, because too many dead ends may make the map unfair or less playable.

During each generation, the generator calculates the fitness score of every map in the current population. Then, maps with better scores are selected as parents through tournament selection. The crossover operation combines partial structures from two parent maps to generate offspring maps. The mutation operation randomly changes a small number of cells with a low probability, such as changing walkable ground into walls or changing walls into walkable ground. After several generations of selection, crossover, and mutation, the system gradually evolves maps that are more suitable for the game environment.

At the end, the map with the highest fitness score is saved as a text file. The map file uses a space-separated \texttt{0} and \texttt{1} format, where each line corresponds to one row of the map. This format is simple, readable, and easy for both Python and C++ game modules to load. Apart from the map grid itself, the system also validates the spawn points of the human and monsters. The spawn points must be located on walkable cells, and the human and monster spawn points cannot overlap. The system also checks the BFS path distance between them. In this implementation, the valid starting distance is set between $10$ and $25$ BFS steps, so that the monster is not placed too close or too far away from the human at the beginning of the game.